\section{A1: the question, the older route, and the recovered proof}
\label{sec:a1}

\subsection{What A1 says}
The archive contains no statement labelled ``A1'': this is the programme's
label. Cote~149, scan~35, distinguishes the full prime-to-characteristic
fundamental group from its abelian quotient; scan~36 asks whether the
point-separation assertion survives in characteristic $p$, and whether its
covering condition forces $P=Q$. The exact displayed maps denoted $(*)$ and
$(**)$ occur before the supplied scans. We do not reconstruct their omitted
wording. The following is the curve statement of the runs, with all
quantifiers restored from the recovered full manuscript
\cite[A.1, A.31--A.33]{A1run}.

\paragraph{Question A1.}
Let $K$ be finitely generated over $\mathbb F_p$ and let $X/K$ be a smooth
geometrically connected hyperbolic curve. Put
\[
 \Delta_X=\pi_1(X_{\overline K})^{(p')},\qquad
 \Pi_X=\pi_1(X)/\ker\bigl(\pi_1(X_{\overline K})\to\Delta_X\bigr).
\]
The arithmetic quotient is the \emph{whole} $G_K$; only the geometric
kernel is replaced. Is
\[
 X(K)\longrightarrow
 \operatorname{Sec}(\Pi_X\to G_K)/\Delta_X,
 \qquad P\longmapsto[s_P]
 \tag{A1}
\]
injective? In characteristic zero the corresponding assertion uses every
geometric prime. Hyperbolic means $2g-2+r>0$ for the geometric smooth
compactification and its $r$ missing points.

For a proper curve of genus at least two over a finite field, the equivalent
covering formulation is this: for every finite constant extension $L/k$
and every connected geometrically prime-to-$p$ Galois cover $D\to C_L$
with a chosen $P'\in D(L)$ above $P$, there is $Q'\in D(L)$ above $Q$
such that $[Q'-P']\in J_D(L)[p^\infty]$. The term ``$p$-torsion'' in
scan~36 has this primary meaning, as its preceding phrase about order a
power of $p$ confirms. The phrase ``a cover of $X$'' in that scan must be
read as a cover of $Y$ in its local notation; see the transcription table.

\paragraph{Finding A01 (Kummer and the covering formulation).}
The preceding covering condition is equivalent to equality of point-section
classes over a finite field. Equality of sections always implies the
stated primary torsion condition. No number-field GRH occurs here.

\begin{proof}
For a finite field $L$, the prime-to-$p$ Kummer kernel of $J_D(L)$ is its
$p$-primary subgroup: finite Kummer maps inject $J_D(L)/nJ_D(L)$ into
$H^1(L,J_D[n])$, and $J_D(L)$ is finite. Equal sections give the same
action on the fibre functor, hence compatible rational lifts to any
pointed cover, with equal abelian Kummer classes.

For the converse we first establish uniqueness of the specified lift.
If two lifts differ by a nonidentity deck element $a$ of order $m$ prime
to $p$, then $p^e(aR-R)$ is principal. Choose a function with this
divisor. Its norm for $\langle a\rangle$ is constant; over the algebraic
closure scale it to norm one. Hilbert~90 writes it as $a(u)/u$.
Thus $\operatorname{div}(u)-p^eR$ is $a$-invariant. The free deck action
has every orbit of length $m$, so the degree $-p^e$ of this divisor is
divisible by $m$, a contradiction.

Take the cofinal characteristic open normal subgroups $N$ of the finitely
generated geometric group $\Delta_C$. The section $s_P$ makes
$N\rtimes s_P(G_k)$ a pointed cover. Make its geometric deck group
constant by a finite extension and apply the assumed condition.
Uniqueness descends the lift of $Q$ to $k$; norm maps and uniqueness
make these lifts compatible as $N$ varies. They form a fixed point of
the profinite path torsor from $P$ to $Q$, which is exactly a geometric
conjugacy of the two sections. The prime-to-$p$ topology on each cover
is the full one: normal cores embed the additional quotient in a finite
product of prime-to-$p$ groups. This also justifies the use of its
Jacobian Kummer map. This checks the form in \cite[A.31]{A1run}.
\end{proof}

\subsection{The exact contribution of Hooley}

\paragraph{Finding A02 (the analytic consequence, conditional on M).}
For every integer $a\ne-1$ which is not a square, write
$K_n(a)=\mathbb Q(\zeta_n,a^{1/n})$. Under M, Hooley's primitive-root
theorem applies, since the trivial Hecke character of $K_n(a)$ has
$L$-function $\zeta_{K_n(a)}$. In particular,
\[
 \#\{\ell\le x:\operatorname{ord}_{\ell}(2)=\ell-1\}
 = A\operatorname{Li}(x)
   +O\!\left(\frac{x\log\log x}{\log^2x}\right),\qquad
 A=\prod_{r\ {
m prime}}\left(1-\frac1{r(r-1)}\right)>0.
 \tag{A02}
\]
This is a cited consequence of \cite[\S7, concluding theorem]{Hooley}; the precise modern
numbered formulation used for the asymptotic and error is
\cite[Theorem~1.1, specialized to $a=2$]{FanPollack}. We do not re-prove
Hooley's sieve. The needed fields have squarefree $n$; already $K_3(2)$
has nonabelian Galois group $S_3$. The stipulated Alp\"oge Dirichlet
theorems therefore do not supply this application of Hooley.

\begin{proof}
Apply M to the trivial character over each $K_n(2)$, and then the
cited theorem. Its density series is
$\sum_{n\ge1}\mu(n)/[K_n(2):\mathbb Q]=A$. The polynomial $T^3-2$
is irreducible by Eisenstein and has nonsquare discriminant $-108$;
its splitting field is $K_3(2)$ and has group $S_3$. This identifies
a specific nonabelian analytic input absent from the Dirichlet scope.
It does not assert that Dirichlet GRH logically contradicts, or can
never imply by some different argument, the needed Dedekind GRH.
\end{proof}

\paragraph{Finding A03 (what the theta route actually needs).}
For a genus-two curve $C/\overline{\mathbb F}_2$, a sufficient geometric
input is the following. The theta divisor is defined by the bundle of
locally exact differentials on the Frobenius twist $C^{(1)}$; we
transport its geometric line bundles to $C$ by the absolute
Frobenius-twist identification, and denote the resulting divisor by
$\Theta_B$. The required condition is
\[
 \begin{gathered}
 \exists\ell>2\text{ prime},\qquad \operatorname{ord}_{\ell}(2)=\ell-1,\\
 \exists L\in\operatorname{Pic}^0(C)[\ell]\setminus\{0\}
                  \quad\text{with }L\in\Theta_B.
 \end{gathered}
 \tag{T(C)}
\]
Together with the necessary covering implication of A01, this implies
point separation for any arithmetic model of $C$. Hooley supplies the
first condition on $\ell$, not the joint assertion $T(C)$.

\begin{proof}
We spell out the geometric implication, so that the exact missing input
is visible. The line bundle $L$ defines a connected cyclic etale cover
$D\to C$ of degree $\ell$. Its nontrivial character spaces
$H^1(C,L^i)$, $1\le i<\ell$, are one-dimensional by Riemann--Roch.
Absolute Frobenius sends the $i$-space to the $2i$-space semilinearly.
The locally exact differential sequence
$0\to\mathcal O_{C^{(1)}}\to F_*\mathcal O_C\to B\to0$,
tensored with $L^i$, identifies its kernel with
$H^0(C^{(1)},B\otimes L^i)$, with the usual Frobenius-twist
identifications. A theta incidence therefore makes one arrow zero.
The primitive-root assumption makes all nontrivial characters one
cycle, so Frobenius on their sum is nilpotent. The Artin--Schreier
sequence identifies its stable part with etale $\mathbb F_2$
cohomology. Thus the new factor $\operatorname{im}(\sigma-1)$ of
$J_D$ has two-rank zero.

Suppose distinct points collide. After a finite extension A01 gives
lifts $P',Q'$ with $T=[Q'-P']$ two-primary. Since
$(\sigma-1)T$ lies in a two-rank-zero abelian variety, it is zero.
Hence $\sigma Q'+P'\sim Q'+\sigma P'$. Equality of the divisors
is impossible by freeness of the deck action and $P'\ne Q'$.
They therefore give a degree-two pencil on $D$. But $D$ cannot be
hyperelliptic: an odd-order free deck automorphism would induce a
nontrivial automorphism of $\mathbb P^1$, which has a fixed point;
its fibre has at most two geometric points, each fixed by the
odd-order element. This contradicts freeness. The implication follows.
The same proof is the prime-degree instance of the older intake's
A.C2.7 and A.C2.19, without requiring its general composite-conductor
criterion.
\end{proof}

\paragraph{Finding A04 (the density inference is false).}
Infinitude, or positive density, of theta-supporting primes would not
by itself combine with Hooley to prove $T(C)$. Indeed the primes
$\ell\equiv1\pmod8$ have density $1/4$, while both $2$ and $-1$
are squares modulo every such prime. Consequently
$\langle2,-1\rangle\ne\mathbb F_\ell^\times$ throughout that set.
This refutes the proposed inference between properties of sets of
primes; it is not a curve counterexample to A1.

\begin{proof}
Quadratic reciprocity gives $(2/\ell)=1$ and $(-1/\ell)=1$ in the
specified progression. Their subgroup lies in the proper subgroup
of squares. Dirichlet's prime-distribution theorem gives density
$1/\varphi(8)=1/4$. The same obstruction persists at $\ell^a$ by
reduction modulo $\ell$. A fixed finite Chebotarev condition can be
combined with primitive roots using \cite[Theorem~3.1, and
Theorems~4.1--4.2 for positivity]{Lenstra1977}; no such encoding,
or positivity for the curve-dependent torsion condition, is supplied
by M. Goldfeld's average theorem \cite[Theorem, p.~223]{Goldfeld}
averages the integer base, and does not supply it for the fixed base~2.
\end{proof}

There is a second limitation in higher genus. The spaces $H^1(C,L^i)$
then have dimension $g-1$. One singular arrow round a Frobenius orbit
need not make its composite nilpotent: take every arrow to be the
identity except one $\operatorname{diag}(0,1,\ldots,1)$.
Its stable subspace has dimension $g-2$. Thus replacing ``ordinary''
by ``not ordinary'' cannot supply zero primitive $p$-rank. The prime-to-$p$ cover statements in
\cite[Propositions~2.3--2.4]{PopSaidi} do not repair this
step. In the latter source ``pairwise prime orders'' is justified by
a proof of pairwise \emph{coprime} orders, not of prime integer orders.

\subsection{A1 does have a different proof in the later run}

\paragraph{Finding A05 (A1, independently of A and M).}
The assertion A1 holds, using the full arithmetic quotient. The
following proof is a checked reconstruction of
\cite[A.2--A.5, A.14--A.19, A.30--A.33]{A1run}, recovered from the
exact checkpoint identified by the supplied summary. It is not a
proof through Hooley, and no originality or formal verification is
claimed. Its decisive input is cuspidal transport with integral
Frobenius information, together with the already known finite
exceptional-set theorem of Sa\"idi--Tamagawa.

\subsubsection*{Step 1. The relative configuration fibre is retained}
For the finite-field assertion let $k$ be finite, let $C/k$ be proper
of genus at least two, and let $P,Q\in C(k)$. To prove the required
geometric group comparison first work with a complex proper curve
$C$ of genus $g\ge2$, and write
\[
 \Lambda=\pi_1^{\rm top}(C-P),\quad
 \Gamma=\pi_1^{\rm top}(C^2-\operatorname{diag}),\quad
 B_0=\pi_1^{\rm top}(C),\quad \Delta=B_0^{(p')}.
\]
The second projection is a locally trivial fibre bundle, obtained
locally by moving the missing point with an isotopy. Since
$\pi_2(C)=0$, it gives $1\to\Lambda\to\Gamma\to B_0\to1$.
Put $F=\Lambda^{(p')}$, $E=\Gamma^{(p')}$,
$K_F=\ker(F\to\Delta)$ and
$K_E=\ker(E\to\Delta\times\Delta)$. Fix $r\ne p$ and replace
each $K$ by its maximal pro-$r$ quotient. Denote the resulting
groups by $F_c,E_c$. We prove
\[
 1\longrightarrow F_c\longrightarrow E_c\longrightarrow\Delta
 \longrightarrow1. \tag{A05.1}
\]

We need a free-group lemma. If $1\to L\to H\to B\to1$ has $L$
free of finite rank at least two, and $H$ acts on
$L/[L,L]L^r$ through an $r$-group, then $L^{(r)}\to H^{(r)}$
is injective. For a finite $r$-group $Q$, the kernel of
$\operatorname{Aut}(Q)\to\operatorname{Aut}(Q/\Phi(Q))$ is an
$r$-group: it acts freely on the $|\Phi(Q)|^d$ ordered generating
lifts of a fixed basis. Characteristic Frattini quotients therefore
show that the closure of the action on $L^{(r)}$ is pro-$r$.
It extends to $H^{(r)}$ and restricts on $L^{(r)}$ to the inner
action. This action is faithful, because a nonabelian free pro-$r$
group has trivial centre. One direct verification uses its Magnus
embedding into $1+(X_1,\ldots,X_d)$ in the noncommutative formal
power-series algebra over $\mathbb F_r$. Injectivity follows by
substitution into each finite $r$-group algebra, whose augmentation
ideal is nilpotent. A series commuting with both $X_1$ and $X_2$
is scalar, so a central Magnus unit of augmentation one is one.
Compactness makes the injection a closed embedding; hence every
finite $r$-quotient of $L$ is dominated by one of $H$.

Now take any finite quotient $Q$ of $F_c$, with normal $r$-subgroup
$R$ the image of $K_F$, and put $A=Q/R$. The corresponding proper
Galois cover $C'\to C$ gives
$Y=(C'\times C')-(a\times a)^{-1}(\operatorname{diag})$.
The removed components are disjoint labelled graphs. Its fibre
$V'=C'-a^{-1}(P)$ has free fundamental group $L$. In
$H_1(V',\mathbb F_r)$, monodromy is trivial both on boundary
classes and on the quotient $H_1(C',\mathbb F_r)$, since the
compactified family is a product. It is therefore an $r$-group of
upper block unipotent transformations. The lemma supplies a finite
$r$-quotient of $H=\pi_1^{\rm top}(Y)$ dominating the restriction
of $Q$ to $L$. Taking its core under the finite quotient
$\Gamma/H=A\times A$ gives a quotient of $\Gamma$ which is an
$r$-group extended by $A\times A$, still prime to $p$.
Its image of $K_E$ is an $r$-group, so it factors through $E_c$.
On the fibre its kernel is contained in $\ker(F_c\to Q)$.
Every finite quotient of $F_c$ is thus detected in $E_c$, proving
injectivity. Right exactness follows from the universal property of
pro-prime-to-$p$ completion; the further killed subgroup lies in
the fibre kernel. This proves (A05.1).

Lift the smooth proper curve in characteristic $p$ to characteristic
zero (SGA~1, III, Corollary~7.4); its marked points lift by smoothness
over the complete henselian base. Apply the
prime-to-$p$ comparison for smooth proper pairs to the pointed
curve and to its square with diagonal boundary (SGA~1, XIII,
Corollary~2.9, \S2.10 and Lemma~2.11). Functoriality identifies
the groups, the relative kernels, and boundary inertia. Thus
(A05.1) holds for the etale groups. Arithmetic fibre comparison
over a finite field is an isomorphism on geometric kernels and
on the unchanged $G_k$, hence on the whole extensions.

\subsubsection*{Step 2. Equal proper sections transport the cusp}
In the arithmetic configuration quotient the diagonal boundary has
a common decomposition group $\mathcal D$ with
\[
 1\to\mathbb Z_r(1)\to\mathcal D\to\Pi_C\to1,
 \tag{A05.2}
\]
mapping diagonally to $\Pi_C\times_{G_k}\Pi_C$. Here is why no
boundary injection is assumed without proof. The boundary of the
$n$th root construction along the diagonal is the gerbe of roots
of its normal line bundle. Finite tame covers extend on its
Kummer charts by SGA~1, XIII, Proposition~5.2. The inverse-limit
boundary group therefore has kernel the image of tame inertia.
After the relative quotient, this is a quotient of $\mathbb Z_r(1)$.
It is the full group: the boundary word of the punctured fibre is
the product of handle commutators, and sending one handle pair
to $1+e_{12},1+e_{23}$ in $\operatorname{UT}_3(\mathbb Z_r)$
sends it to $1+e_{13}$ of infinite order. Step~1 retains this
inertia. The transverse fibre has multiplicity one, so its cusp
group is the full inverse image of $s_P(G_k)$ in $\mathcal D$.

If $s_P,s_Q$ are geometrically conjugate, lift their conjugator
$\gamma\in\Delta_C$ through $\mathcal D$. Its conjugation
identifies the two fibre groups and cusps, and acts on their first
proper coordinate by $\operatorname{Inn}(\gamma)$. A geometric
element of the target fibre mapping to $\gamma^{-1}$ cancels this
inner automorphism. Consequently we obtain
\[
 \beta:G_P\xrightarrow\sim G_Q
 \quad\text{over }\operatorname{id}_{\Pi_C},
 \tag{A05.3}
\]
where $G_P,G_Q$ are the maximal cuspidally pro-$r$ quotients of
the actual punctured curves. It preserves cusp conjugacy classes.

\subsubsection*{Step 3. Corresponding open covers have the same counts}
If $K\triangleleft F$ is closed and
$N=\ker(K\to K^{(r)})$, every open $I\le F$ containing $N$
satisfies $I^{(r)}\simeq(I/N)^{(r)}$. Indeed take a finite
$r$-quotient $I/J$, put $I_0=\operatorname{core}_F(I)$, and take
the core $M$ of $J\cap I_0$. Then $I_0/M$ is an $r$-group.
The image of $K$ in $F/I_0$ is also an $r$-group because
$N\subset I_0$. Thus its image in $F/M$ is an $r$-group, and
the defining property of $N$ gives $N\subset M\subset J$.

Apply this to a geometric open cover of $G_P$ or $G_Q$. Killing
its actual cusp inertias and taking its pro-$r$ quotient recovers
the full pro-$r$ group of its smooth compactification. This is
equivariant for arithmetic Frobenius. Thus $\beta$ identifies
the integral $r$-adic $H^1$ of corresponding compactifications.
The trace formula gives identical compactified point counts over
every finite extension; the cusp double cosets give identical
degrees of the missing points. Subtracting them gives identical
counts for the open covers. The inputs are the usual curve trace
formula and purity, not analytic RH over $\mathbb Q$.

\subsubsection*{Step 4. Compactness recognizes actual point sections}
Fix $R\in(C-P)(L)$ and let $S=\beta(D_R)$. For every open
subgroup $H$ of $(G_Q)_L$ containing $S$, the source cover for
$\beta^{-1}H$ has an $L$-point above $R$. Step~3 shows that the
corresponding target cover $X_H$ has an $L$-point. Its finite
sets $X_H(L)$ have nonempty inverse limit: any finite collection
of compatibility conditions is satisfied by a point on the cover
for the intersection of the corresponding subgroups. Compactness
then supplies a compatible family, and a lift of its base point
$T\in(C-Q)(L)$ to the full profinite fibre. Its decomposition
group lies in every $H$ (rationality at each level, and injectivity
of the decomposition group's projection to $G_L$), hence in
$\bigcap_{H\supset S}H=S$. Both project isomorphically to $G_L$,
so they are equal. Thus $\beta$ transports every point section
to a genuine point section.

Sa\"idi--Tamagawa \cite[Definitions~1.3, 1.5, 1.10 and
Proposition~1.8(vi)]{ST2009} give one finite exceptional image
$E_C\subset C^{\rm cl}$ for the cofinite prime set $p'$.
Outside it the proper decomposition group uniquely determines
the geometric point, even up to commensurability. Remark~1.2
makes this one exceptional set valid over every constant extension.
Since $\beta$ is over the identity of $\Pi_C$, the transported
point is $R$ whenever $R\notin E_C\cup\{P,Q\}$.

\subsubsection*{Step 5. The Kummer difference is impossible}
The proper abelianization gives $[P-Q]\in J_C(k)[p^\infty]$.
Choose $f\in k(C)^\times$ with
$\operatorname{div}(f)=p^e(P-Q)$; descent of a geometric
generator follows from Hilbert~90. Put $M=\mathbb Z_r(1)$.
The geometrically cuspidally central quotients $E_P,E_Q$ have
kernel $M$. Infinite inertia follows from the same Heisenberg
test, and their geometric extension classes are the degree-one
point classes in $H^2(\Delta_C,M)=\mathbb Z_r$.
The map induced by $\beta$ acts on $M$ by a unit $u$.
Since it fixes the proper quotient and the degree-one class,
$u=1$.

For completeness the orientation identification uses the
punctured surface presentation (SGA~1, XIII, Corollary~2.12).
Equivalently it is the Kummer first Chern class of
$\mathcal O_C(P)$, with residue $+1$. The group/etale
cohomology comparison follows from the tower of all geometric
prime-to-$p$ covers: cyclic torsors kill $H^1$ classes, covers
of degree divisible by $r^a$ kill $H^2(-,\mu_{r^a})$, and a
smooth proper curve has no higher such cohomology. Taking the
associated spectral sequence and then inverse limits gives
the asserted $\mathbb Z_r(1)$ identification.

On $E_P\times_{\Pi_C}E_Q$ the difference
$\chi(u,v)=\overline\beta(u)v^{-1}$ is an $M$-valued
crossed homomorphism. Pull it back to the group of
$C-\{P,Q\}$. Its residues are $(1,-1)$, so
$\kappa_r(f)-p^e\chi$ has zero residues and descends to
$H^1(\Pi_C,M)$. Inflation-restriction identifies this group
with $H^1(G_k,M)$: the Frobenius eigenvalues on
$\operatorname{Hom}(T_rJ_C,\mathbb Z_r(1))\otimes\mathbb Q_r$
have absolute value $\sqrt{\#k}$, hence no invariant, and its
integral module is torsion free. Finite-field Kummer theory
therefore gives $a\in k^\times$ with
\[
 \kappa_r(f)=p^e\chi+\kappa_r(a). \tag{A05.4}
\]

At a good closed point $R$ put $L=k(R)$. Step~4 makes
$\beta(D_R^P)$ conjugate to $D_R^Q$. Its conjugator projects
to the normalizer in $\Pi_{C_L}$ of the proper section $D_R$, which is $D_R$
itself by \cite[Lemma~1.7]{ST2009}. Multiplying by an element
of $D_R^Q$ puts it in the cuspidal kernel; in $E_Q$ it lies
in $M$. Thus $\chi$ evaluates to a coboundary, and (A05.4)
gives $\kappa_{r,L}(f(R)/a)=0$.

If $P\ne Q$, the function $f$ is nonconstant. The excluded
points and their $f$-values form a finite set. Choose a
nontrivial $r$-power root of unity $\zeta\in\overline k$
with $a\zeta$ outside that set, and a point $R$ above it.
Then $\zeta=f(R)/a\in k(R)^\times$. The $r$-adic Kummer
map of this finite cyclic group is injective on its $r$-primary
subgroup, contradicting its asserted zero value. This proves
injectivity over every finite field, including every constant
extension.

\subsubsection*{Step 6. Finitely generated fields and open curves}
For a finitely generated field $K/\mathbb F_p$, spread a proper
curve and two distinct colliding points to a smooth connected
finite-type base, shrinking so that the sections are disjoint.
At a closed point use its henselian local ring $R$ and fraction
field $L$. Smooth proper specialization (SGA~1, X,
Theorem~3.8 and Corollary~3.9) identifies the geometric
prime-to-$p$ groups. The arithmetic map $G_L\to G_{k(s)}$
is onto, since finite etale covers of the normal henselian
domain have connected generic fibres. The same conjugator
would therefore give colliding distinct points on the finite
residue-field curve, a contradiction.

We record also the Kummer fact needed for open curves. For a
semiabelian $G$ over $K$ the prime-to-$p$ Kummer kernel is
$G(K)[p^\infty]$; in characteristic zero, using all primes,
it is zero. After a finite splitting extension of its torus,
the divisor sequence on a normal proper model makes
$K^\times$ an extension of constant units by a free abelian
group. Constant units are finite over finite constants, and
finite torsion plus free abelian over number-field constants.
Lang--Neron (\cite[Theorem~2.1]{Conrad}, with the perfect constant
field algebraically closed in $K$), together with Mordell--Weil on the trace over
number fields, makes the abelian quotient's rational points
finitely generated. Lifting its finitely many generators and
relations shows that $G(K)$ is free abelian plus a finitely
generated group; the intersections of its prime-to-$p$
multiples and of all its multiples are respectively its
$p$-primary torsion and zero. This is the argument of
\cite[A.30]{A1run}; it does not falsely assume $K^\times$
finitely generated.

Make a finite algebraic extension to split the compactification
and boundary; purely inseparable base change leaves the etale
category and Galois group unchanged (SGA~1, IX,
Theorem~4.10). If the compactified genus is at least two,
map a collision to the proper curve. In genus zero, choose
three boundary points as $0,1,\infty$; their coordinate embeds
the curve into $\mathbb G_m$, and equal Kummer classes have
ratio of $p$-power order in $K^\times$, hence ratio one.
In genus one use a rational point as origin of the compact
elliptic curve $E$. A collision difference has order $p^a$.
Pull back the open curve under $[m]$, with $(m,p)=1$ and
$m^2>p^a$, and split its boundary and deck group by finite
extension. Lift the collision with first point zero. Its class
in the generalized Jacobian has exact order $p^a$: the torus
kernel has no nonzero geometric $p$-primary points. Thus a
function with divisor $p^a(Q'-P')$ equals one at all points
of the pulled-back reduced boundary. There are at least
$m^2$ of them, whereas $f-1$ has pole degree $p^a$, impossible.
In characteristic zero the all-prime Kummer kernel is zero
and the Abel map to the generalized Jacobian separates points.
This proves A1 with the stated scope. It does not assert the
full higher-dimensional or section-surjectivity conjecture.

\paragraph{Consequence for the two inputs.}
A05 needs neither A nor M. Therefore it remains true if only
Alp\"oge's stated theorems are added to the existing background.
Hooley's number-field application is still conditional on M;
the theta-support and primitive-cover questions are separate
geometric assertions. The older intake's unresolved A1 status
has been superseded by the recovered later proof, not by an
unproved implication from primitive-root density.
